\section{Introduction}
\label{sec:introduction}

Let $f\colon X\to Z$ be a projective contraction of normal varieties,
meaning that $f$ is surjective and $f_*\OO_X=\OO_Z$. A complement near a
closed point $z\in Z$ is an effective rational boundary $B$ for which the
adjoint divisor $K_X+B$ becomes torsion on the inverse image of a
neighbourhood of $z$. More precisely, an index $N$ records that
$N(K_X+B)$ is integral Cartier and linearly trivial there. The complement
is klt if every log discrepancy of $(X,B)$ is positive. We ask whether one
index works for all contractions of a fixed dimension whose singularities
have a fixed positive lower discrepancy bound.

When $K_X$ is rational Cartier and $-K_X$ is ample over $Z$, we call $f$ a
\emph{Fano contraction}. Complements turn this positivity into an adjoint
divisor with controlled torsion. The distinction between lc and klt
complements is essential: an lc complement permits zero discrepancies,
whereas a klt complement must keep every discrepancy strictly positive.
We prove that a uniform positive bound for the singularities of $X$
suffices to achieve the latter conclusion with a bounded index.

\begin{theorem}\label{thm:main}
For every positive integer $d$ and positive rational number $\epsilon$ there
is a positive integer $N=N(d,\epsilon)$ with the following property. Let $k$
be an algebraically closed field of characteristic zero, and let
$f\colon X\to Z$ be a projective surjective morphism of normal integral
$k$-varieties such that $f_*\OO_X=\OO_Z$, $\dim X=d$, $K_X$ is
$\Q$-Cartier, $X$ is $\epsilon$-lc, and $-K_X$ is ample over $Z$.
For every closed point $z\in Z$ there are an effective $\Q$-divisor $B$ on
$X$ and a Zariski open neighbourhood $U$ of $z$ such that $(X,B)$ is klt over
$U$ and
\[
 N(K_X+B)|_{f^{-1}(U)}\sim 0
\]
is an integral Cartier divisor.
\end{theorem}

The morphism in Theorem~\ref{thm:main} need not have relative Picard number
one, nor need its general fibre have positive dimension. The displayed
linear equivalence means triviality of the associated line bundle on
$f^{-1}(U)$, after the indicated shrinking; it is stronger than numerical
triviality. Both $B$ and $U$ may depend on $f$ and $z$, whereas $N$ does not.

A variety is of Fano type over its base if it admits an effective
$\Q$-boundary $\Delta$ for which $(X,\Delta)$ is klt and
$-(K_X+\Delta)$ is relatively ample. Equivalent versions using a nef and big
anti-log-canonical divisor may be converted to this one by a small boundary
perturbation.

The proof uses the equivalence between bounded klt complements and a local
threshold statement established by Chen \cite[Theorem~1.6]{Chen}. The
birational, boundary-free case of that threshold statement remains the main
issue. We formulate it in Proposition~\ref{prop:birational-threshold} and
prove it by contradiction. The role of bounded lc complements is different:
they supply an auxiliary linear system of a fixed degree $n$. Throughout the
proof, this $n$ is distinguished from the final klt complement index $N$.

\subsection{Earlier results and the finite-coefficient conjecture}

Shokurov introduced complements in his work on threefold log flips
\cite[Section~5]{ShokurovFlips} and formulated their boundedness
systematically in his surface theory \cite[Conjecture~1.3]{ShokurovSurfaces}.
Prokhorov--Shokurov developed the passage from global complements to local
ones and the higher-dimensional inductive program
\cite{ProkhorovShokurovFirst,ProkhorovShokurov}.
Shokurov's later treatment includes a bounded lc complement theorem
assuming the existence of a klt real complement
\cite[Theorem~3]{ShokurovComplements}. Here the klt condition concerns
the input real complement; the bounded-index output is lc.
An lc complement may have places of zero discrepancy. Reducing its
boundary coefficients need not preserve the linear triviality of the
adjoint divisor, so keeping the complement klt requires an additional
argument.

Birkar proved boundedness of lc complements for Fano type fibrations with
hyperstandard coefficients \cite[Theorems~1.7 and~1.8]{BirkarComplements},
and boundedness of klt complements for boundary-free Fano fibrations over curves
\cite[Corollary~1.4]{BirkarFibrations}. Chen related bounded klt complements
to uniform local hypersurface thresholds and established the conjecture for
bases of dimension at most two
\cite[Theorem~1.6 and Corollary~1.8]{Chen}.

Theorem~\ref{thm:main} also gives the following consequence in Chen's
complex quasi-projective setting.

\begin{corollary}\label{cor:finite-coefficients}
Fix a positive integer $d$, a real number $\epsilon>0$, and a finite set
$I\subset[0,1]\cap\Q$. There is a positive integer $M=M(d,\epsilon,I)$
with the following property. Let $(X,\Delta)$ be an $\epsilon$-lc pair of
dimension $d$ with coefficients in $I$, over $\C$. Let $f\colon X\to Z$
be a contraction of quasi-projective normal varieties, assume that $X$ is
of Fano type over $Z$, and assume that $-(K_X+\Delta)$ is nef over $Z$.
For every closed $z\in Z$, there is an effective $\Q$-divisor
$\Delta^+\ge\Delta$ such that, near $f^{-1}(z)$, the pair
$(X,\Delta^+)$ is klt and $M(K_X+\Delta^+)$ is Cartier and linearly
trivial after shrinking the base about $z$.
\end{corollary}

The proof, using Chen's equivalence and a crepant ample model, is given
in Section~\ref{sec:reduction}.

Corollary~\ref{cor:finite-coefficients} is the finite-rational-coefficient
form of Shokurov's bounded-klt-complement conjecture stated in
\cite[Conjecture~1.1]{Chen}. The same equivalence yields the local
hypersurface assertions in \cite[Conjectures~1.3 and~1.4]{Chen}.
The conclusion concerns klt complements; the stronger requirement of an
$\epsilon$-lc complement in \cite[Remark~1.2(1)]{Chen} is not asserted here.

\subsection{Methods and related work}

Our first tools come from normalized-volume minimization. Li established
properness estimates \cite{LiProperness}; Blum proved existence of a
minimizer \cite{BlumExistence}; Xu proved its quasi-monomiality
\cite{XuQuasimonomial}; and Xu--Zhuang proved uniqueness and finite
generation \cite{XuZhuangUniqueness,XuZhuangStable}. Li--Xu identifies
the minimizing Reeb valuation on the resulting K-semistable cone
\cite{LiXuStable}. We use this ordinary klt theory after proving an
additional exactness statement: one fixed small rational perturbation
has its actual global minimizer on the prescribed lc constraint.
This is Proposition~\ref{a:exact-penalization}; its constants are allowed
to depend on the individual germ and ideals.

The projective stage combines minimal models and Mori dream spaces
\cite{BCHM,HuKeel}, Ambro's canonical bundle formula
\cite{AmbroBoundary,AmbroModuli}, and Birkar's boundedness and
threshold theorems \cite{BirkarBAB,BirkarPolarised,BirkarFibrations}.
The needed uniformity concerns both divisor classes and effective
representatives. We retain those representatives in marked bounded
families before passing to generic fibres or taking limits.

The final stage uses Takagi's comparison between multiplier and test ideals
\cite{Takagi}, in the dense-open formulation recalled by
de Fernex--Docampo--Takagi--Tucker \cite[Theorem~11]{DFDTT}, together
with Schwede's Cartier-algebra formulation
\cite{SchwedeNonQG}. These results supply trace generation for a fixed
rational pair. The proof below additionally retains the finite output
spaces and their valuation bounds through two successive liftings.
Related local Cartier-threshold bounds under a positive normalized-volume
lower bound appear in \cite[Theorem~1.7]{HanLiuQiLocalVolumes}.

The companion article \cite[Theorem~1]{OpenAICartier} proves a uniform
Cartier-section threshold for rational-boundary Fano type contractions
by an independent character and pinning argument. Its endpoint is a
Cartier divisor through a prescribed base point with a uniformly lc
pullback. The proof here develops the complement index directly through
the birational threshold and the trace estimates above.

\subsection{Outline of the proof}

The reduction in Section~\ref{sec:reduction} asks for a uniform lower
bound on the log canonical threshold of the pullback of a closed base
point's maximal ideal in the birational, zero-boundary case. A sequence
violating this bound has divisors whose
orders on that ideal are arbitrarily large compared with their discrepancies.
We will construct a prime divisor of discrepancy less than $\epsilon$.
Small discrepancy at a real valuation is not by itself enough, since
rescaling the valuation also rescales its discrepancy. The proof must
therefore control both discrepancy and the eventual divisor normalization.

\smallskip\noindent
\emph{From concentration to a graded model.}
After shrinking each base to an affine neighbourhood of $z$ and replacing
the sequence by its crepant ample models,
Sections~\ref{a:cone}--\ref{a:essential} encode anticanonical sections in
an affine cone. Its bounded-degree lc system has a centred lc place
detecting the concentration of the base ideal. We minimize normalized
volume subject to this lc equality. Proposition~\ref{a:exact-penalization}
shows that the constrained minimizer is exactly the ordinary minimizer of
a small effective klt perturbation. This permits ordinary stable
degeneration to produce a finitely generated graded ring, while a
comparison at the minimizer preserves concentration of the base ideal.

Compatible further filtrations let us maximize the rank of the torus
whose characters record the occurring valuation degrees. Each successive
valuation remains realized on the original cone. The resulting projective quotient carries a
klt pair $(S,B)$ and an effective ample rational divisor
$H\sim_{\Q}-(K_S+B)$ for which $(S,B+H)$ is lc. At each lc place of
this latter pair, replacing $H$ by any effective rationally linearly
equivalent divisor increases its order by at most a factor $1+\eta$,
where $\eta\to0$ along the sequence. This is the weak
essentiality property of Proposition~\ref{a:weak-essentiality}.

We also keep actual sections on $X$. Choose a general degree-$n$
anticanonical section $h$, and let $B_h$ be one $n$-th of its effective
section divisor, so $(X,B_h)$ is lc. For $m\in n\Z_{\ge0}$, dividing
degree-$m$ anticanonical sections by $h^{m/n}$ gives a space of rational
functions $V_m\subset k(X)$. The normalized cone valuation restricts to
a valuation $v$ with $A_{X,B_h}(v)=0$. The cone valuation assigns to a
section the \emph{cost} $b(F)=m+v(F)$ of the corresponding
$F\in V_m$; the assigned
degree $m$ is retained in this notation. Costs are nonnegative,
$V_0=k[Z]$, and $v(\m_z)\to\infty$ along the sequence.
These systems and the exact
discrepancy identity below are constructed in
Subsection~\ref{a:angular}.

\smallskip\noindent
\emph{Upper counts and independent sections.}
Sections~\ref{b:discrepancy-push}--\ref{b:bounded-axes} use weak
essentiality to control these section spaces. A finite tower of projective
contractions removes one small scale of the residual field $k(S)$ at a time.
Horizontal discrepancy bounds give divisor-order estimates on each
relative field. Together with the torus directions, these fields determine
$d$ widths $W_j$: independent rational functions in direction $j$ can be
obtained at degree and cost $O(W_j)$.

The two needed estimates go in opposite directions. At any fixed physical
degree, the sum of initial spaces through cost $T$ has dimension at most
a uniform multiple of $\prod_j(1+T/W_j)$. Monomials in actual sections
with independent initials give lower box bounds at comparable degree and
cost. Here an
initial is the least-valued part of a function in the associated graded
field. The directions whose widths stay bounded need a further estimate:
through a slowly growing cutoff, cost is bounded by a uniform multiple
of physical degree, and
the dimension is $O((1+m)^r)$, where $r$ is their number. To prove this,
Section~\ref{b:bounded-axes} retains effective divisor representatives in
a marked projective family. An obstruction would produce a degree-zero
section of bounded positive cost, contradicting concentration of the
entire base maximal ideal.

\smallskip\noindent
\emph{A small discrepancy with controlled divisor normalization.}
The connection back to the singularities of $X$ is
\[
 A_X(w)=A_{X,B_h}(w)+
       \sup_{\substack{m\in n\Z_{>0}\\0\ne F\in V_m}}
                          \frac{-w(F)}m
\]
for quasi-monomial valuations centred on $X$; this is \eqref{eq:A13}.
Thus inequalities
$w(F)+B_*m\ge q\,b(F)$, with $B_*,q>0$, bound the second term by
$B_*$. In degree zero they give nonnegative values on $k[Z]$ and
positive values on $\m_z$, forcing the centre to lie over $z$.
Sections~\ref{c:optimization-section}--\ref{c:rank-section} seek such
valuations with small $A_{X,B_h}(w)$, while preserving independent
initial coordinates. The fractional-form discrepancy $A_\theta$ used
there is precisely $A_{X,B_h}$.

The optimization has a finite test and compact sublevels on a monomial
skeleton. A suitably chosen optimizer gives a strict divisorial budget on a
geometric generic fibre of the full initial field, which retains all
homogeneous degrees, rather than just residue classes. In suitable
positive-characteristic reductions, Cartier trace converts this budget
into sections of controlled degree and value.
Relative trace survival is first lifted to the full initial field and
then to the original field; the finite output spaces are fixed before
reduction to positive characteristic.

Comparing these trace outputs with the upper counts controls both changes
of coordinate values and growth of the number of independent initials.
First the bounded-width coordinates acquire value zero. A further trace
argument then gives two-sided bounds for section values. Through the
remaining width layers, independent initials are added and their values
are rounded to integers, preserving all earlier constraints. Once the
initials form a full transcendence basis, the section counts bound the
denominator of the entire value group. A bounded multiple of the resulting
valuation is an integral multiple of a prime divisor. The final choice of
tolerances and scaling makes that divisor's discrepancy less than
$\epsilon$, completing the contradiction.

\subsection{Conventions}

All varieties in characteristic zero are over an algebraically closed field.
After Section~\ref{sec:reduction} we work over $\C$ through
Section~\ref{b:bounded-axes}. In Section~\ref{c:optimization-section} we
descend the contradiction sequence and its specified data to a countable
algebraically closed field of characteristic zero; later constructions allow
the stated extensions of constants. From Section~\ref{c:trace-section},
positive-characteristic reductions are auxiliary trace calculations on
retained finite data: the chosen valuations, optimization, and strict
discrepancy assertions remain in characteristic zero. A contraction
is a projective surjective morphism with connected fibres in the sense
$f_*\OO_X=\OO_Z$.

We use log discrepancies. For a prime divisor $E$ on a smooth birational model
$\pi\colon Y\to X$,
\[
 A_{X,B}(E)=1+\ord_E\bigl(K_Y-\pi^*(K_X+B)\bigr).
\]
Discrepancies, divisor orders, and ideal orders are extended homogeneously to
positive real multiples of valuations. The order of an ideal is the minimum
of the orders of local generators. The notation $A_X$ means $A_{X,0}$.
For a nonzero ideal $\mathfrak a$ on a klt variety,
\[
 \lct(X;\mathfrak a)=
 \inf_{E(\mathfrak a)>0}\frac{A_X(E)}{E(\mathfrak a)},
\]
where the infimum is over prime divisors over $X$. Rational Cartier divisors
are evaluated by pullback on a model on which the valuation has a centre.

In contradiction sequences we pass to subsequences and shrink the base about
the chosen point without repeating this convention. Constants called uniform
are independent of the sequence index after these operations; their allowed
geometric dependencies will be stated. The number of contractions, scale
layers, and rank increases is bounded. Small positive constants chosen before
a limit are kept separate from quantities tending to zero along the sequence.
The original dimension is $d$; the first section cone has dimension $p=d+1$.
Characteristic primes in the trace arguments are denoted by $\mathfrak p$.
